\documentclass[10pt,a4paper]{amsart}
\usepackage[margin=19mm]{geometry}
\usepackage{amsmath,amssymb,mathtools}
\usepackage[scr=rsfs,cal=euler]{mathalfa}
\usepackage{xfrac}
\usepackage[unicode,hidelinks,bookmarks=false]{hyperref}
\newcommand{\ie}{i.e.\ }
\newcommand{\cf}{cf.\ }
\newcommand{\resp}{respectively\ }
\newcommand{\Subsection}{\subsection}
\providecommand{\nobreakdash}{\nobreak\hbox{-}}
\setlength{\emergencystretch}{2em}
\multlinegap0pt
\usepackage{amssymb}
\newtheorem{thm}{Theorem}
\title{Subgroups of Finite Fields As Cap Sets}
\author{Anthony Kable, Melissa Mills, David J. Wright}
\date{Source: arXiv:2604.26989v1 (CC BY 4.0)}
\begin{document}
\maketitle

\noindent\emph{Source and licence.} Subgroups of Finite Fields As Cap Sets. Version arXiv:2604.26989v1 (CC BY 4.0), distributed by arXiv under the Creative Commons Attribution 4.0 International licence, http://creativecommons.org/licenses/by/4.0/, as recorded in the arXiv licence metadata for this exact version and shown on the abstract page https://arxiv.org/abs/2604.26989v1. Full licence notice: \texttt{licenses/arxiv-2604.26989-CC-BY-4.0.txt}.\par\smallskip

\noindent\textit{Editorial scope.} Counted unit: the article's thm thm (source0.tex lines 252--262). The article's complete original proof of it is delivered with it (source0.tex lines 264--279, one \texttt{proof} environment of 16 lines). No mathematics is written, repaired or re-derived: the statement and the proof are the article's own lines, in the article's own notation, and no retained line is substituted or re-flowed.  No further source-local context is needed: every cross-reference of every retained line resolves inside the two delivered spans.  Nothing was omitted from any retained span, so the package carries no omission record.  No retained line contains a citation, so the wrapper carries no bibliography and no bibliography entry.  No retained line contains a figure environment, an image-inclusion command or drawing code, so no image is delivered and the package declares no asset dependency.  Editorial formatting only, in addition: an AMS \texttt{amsart} preamble that restates the article's own declarations and macros used by the retained lines, namely the environments \texttt{thm} on separate counters, together with the standard packages those lines need.  The retained environments print in this document as Theorem 1 in order of appearance, whereas the article prints the same items with the numbering of its own full document.  The complete native source of the article as distributed in the arXiv e-print for this exact version is bundled unchanged as \texttt{sources/199438/source0.tex}.
\begin{thm}
  For a prime power $q$ and divisor $d\mid q-1$, suppose the subgroup
  $G$ of $d$-th powers in $\mathbb{F}_q^\times$ is a cap in the sense
  that there are no solutions to $x_1+\dots+x_m=0$ with distinct
  $x_1,\dots,x_m\in G$. An element $y\notin G$ is represented by
  $G$ if there are distinct $x_1,\dots,x_{m-1}\in G$ with
  $y+x_1+\dotsb+x_{m-1}=0$. We say $G$ is complete if every
  $y\notin G$ is represented by $G$.  Let $g$ be a primitive root of
  $\mathbb{F}_q$.  Then $G$ is complete if 0 and all powers $g^k$ for
  $1\le k\le d-1$ are represented by $G$.
\end{thm}

\begin{proof}
  We must prove every $x\notin G$ is represented by $G$. By
  hypothesis, we may assume $x=g^k$ for some integer $d\le k\le q-2$
  and $d\nmid k$. By division, $k=ad+r$ for integers $a$ and $r$ with
  $1\le r\le d-1$. Since we assumed $g^r$ is represented by $G$, we
  have:
  $$
  g^{r} + g^{b_1d}+\dotsb+g^{b_{m-1}d}=0,
  $$
  for integers $b_1,\dots,b_{m-1}$. Multiplying by $g^{ad}$ shows
  $g^k=g^{ad+r}$ is represented by $G$:
  $$
  g^{ad+r} + g^{(b_1+a)d}+\dotsb+g^{(b_{m-1}+a)d}=0.
  $$
  Thus,   $G$ is complete. 
\end{proof}

\end{document}
